% !PRINTED 615
\weberpage{615}

%<<CONT
or, if one carries out the composition,
\weq{7}{\varphi(z + 1) \equiv a' \varphi(z) + a\ (\mathrm{mod}\ n)}
for every integer $z$. Since $\varphi(z)$ does not reach the degree $n$, the coefficients of the individual powers of $z$ on the two sides of (7) must (\S.\,136) be congruent, and from the comparison of the coefficients of the highest powers of $z$ it follows that $a' \equiv 1\ (\mathrm{mod}\ n)$, and hence
\[
\varphi(z + 1) \equiv \varphi(z) + a\ (\mathrm{mod}\ n).
\]
If one here puts $z + 1, z + 2 \ldots z + h - 1$ for $z$, it follows that
\[
\varphi(z + h) = \varphi(z) + ah,
\]
where $h$ can be congruent to any arbitrary integer with respect to the modulus $n$. In this one may now set $z = 0$, and obtains, on writing $z$ again for $h$ and setting $\varphi(0) = b$,
\weq{8}{\varphi(z) = az + b.}

The group $P$ therefore consists entirely of linear permutations, and our theorem is proved.

We add further the remark, which follows immediately from an inspection of the formulae, that the \gesperrt{cyclic group is a normal divisor of every linear group in which it is contained at all}.

If we apply the theorem just proved to the group of the metacyclic equation, we obtain the theorem of \gesperrt{Galois}:

II. \gesperrt{The group of a metacyclic equation of prime degree is linear.}

For if we return to the chain of groups (1), we have seen that $P_{\mu - 1}$ is the cyclic group. If $P_{\mu - 1}$ is not identical with $P$, then $P_{\mu - 1}$ is a normal divisor of $P_{\mu - 2}$, and hence $P_{\mu - 2}$ is linear. $P_{\mu - 2}$ is in turn a normal divisor of $P_{\mu - 3}$, hence $P_{\mu - 3}$ too is linear, and so on, until we reach the conclusion that $P$ itself must also be linear. At the same time it follows that $P_{\mu - 1}$ is a normal divisor of all the preceding groups, and hence also of $P$ itself.

The converse theorem also holds:

III. \gesperrt{Every irreducible equation of prime degree whose group is linear is metacyclic.}

To prove it, it suffices to show that every transitive linear group $L$ has a normal divisor $L'$ whose index
%CONT>>
